% ============================================================================
% ESQUEMA PROPOSTO (banca, sugestão — não mergeado): o desenho da avaliação
% da composição do L0 e da partida a frio (P1/P2 no vocabulário da defesa):
% três origens para o conjunto inicial, um único protocolo de medição, e a
% nota de anti-circularidade do envelope. Ilustra a Seção sec:metodo-l0 e
% complementa o esq-drisl (que mostra o algoritmo; este mostra a avaliação).
%
% Uso no Cap. 3 (dentro de um ambiente figure):
%   % ============================================================================
% ESQUEMA PROPOSTO (banca, sugestão — não mergeado): o desenho da avaliação
% da composição do L0 e da partida a frio (P1/P2 no vocabulário da defesa):
% três origens para o conjunto inicial, um único protocolo de medição, e a
% nota de anti-circularidade do envelope. Ilustra a Seção sec:metodo-l0 e
% complementa o esq-drisl (que mostra o algoritmo; este mostra a avaliação).
%
% Uso no Cap. 3 (dentro de um ambiente figure):
%   % ============================================================================
% ESQUEMA PROPOSTO (banca, sugestão — não mergeado): o desenho da avaliação
% da composição do L0 e da partida a frio (P1/P2 no vocabulário da defesa):
% três origens para o conjunto inicial, um único protocolo de medição, e a
% nota de anti-circularidade do envelope. Ilustra a Seção sec:metodo-l0 e
% complementa o esq-drisl (que mostra o algoritmo; este mostra a avaliação).
%
% Uso no Cap. 3 (dentro de um ambiente figure):
%   % ============================================================================
% ESQUEMA PROPOSTO (banca, sugestão — não mergeado): o desenho da avaliação
% da composição do L0 e da partida a frio (P1/P2 no vocabulário da defesa):
% três origens para o conjunto inicial, um único protocolo de medição, e a
% nota de anti-circularidade do envelope. Ilustra a Seção sec:metodo-l0 e
% complementa o esq-drisl (que mostra o algoritmo; este mostra a avaliação).
%
% Uso no Cap. 3 (dentro de um ambiente figure):
%   \input{3-metodo/esquemas-propostos/esq-l0-tres-origens}
% Uso standalone (pré-visualização):
%   \documentclass[tikz,border=6pt]{standalone}
%   \usetikzlibrary{arrows.meta, positioning}
%   \begin{document}\input{esq-l0-tres-origens.tex}\end{document}
%   (no modo standalone, troque as \ref{...} por texto fixo ou defina-as)
% Requer apenas arrows.meta e positioning (já em packages.tex). Legível em
% P&B: origem = branco; protocolo = cinza claro; leitura = cinza médio.
% FREEZE respeitado: todos os números são os da própria Seção 3.6
% (47 tamanhos, 10 a 200.000, 30 repetições; 4 cenários, 10 tamanhos,
% 10 a 30.000), nenhum resultado afirmado — a leitura remete ao Cap. 4.
% ============================================================================
\begin{tikzpicture}[
    x=1cm, y=1cm,
    caixa/.style={draw, rounded corners=2pt, align=center, text width=42mm,
                  minimum height=11mm, inner sep=3pt, font=\footnotesize},
    protocolo/.style={caixa, fill=gray!8},
    leitura/.style={caixa, fill=gray!14},
    seta/.style={-{Stealth[length=2.4mm]}, thick},
    rot/.style={font=\scriptsize, align=center, inner sep=2pt}
]
    % ------------- três origens para o conjunto inicial --------------------
    % (o AG fica na linha de baixo, vizinho da sua nota de anti-circularidade:
    % com ele no meio, a seta tracejada atravessava a caixa do DRI-SL)
    \node[caixa, text width=44mm] (aleat) at (0.6,0)
        {\textbf{amostragem aleatória}\\(dispersão do sorteio):\\
         47 tamanhos de $L_0$, de 10 a 200.000; 30 repetições por tamanho};
    \node[caixa, text width=44mm] (drisl) at (0.6,-2.9)
        {\textbf{DRI-SL} (proposto):\\determinístico e sem acesso a
         rótulos; algoritmo na Seção~\ref{sec:metodo-drisl}};
    \node[caixa, text width=44mm] (ag) at (0.6,-5.8)
        {\textbf{algoritmo genético}\\(envelope: teto e piso):\\
         4 cenários (max/min de Acurácia\\e de Macro F1) em 10 tamanhos,\\
         de 10 a 30.000};

    % ------------- protocolo único e leitura -------------------------------
    \node[protocolo, text width=46mm] (proto) at (6.8,-2.9)
        {protocolo comum a toda origem:\\o classificador PVBin é treinado\\
         no $L_0$; Acurácia e Macro F1\\medidas no conjunto de teste};
    \node[leitura, text width=46mm] (leit) at (6.8,-6.4)
        {leitura nos resultados:\\
         sensibilidade (Seção~\ref{sec:res-l0-sens}),\\
         envelope (Seção~\ref{sec:res-l0-ag}),\\
         \mbox{DRI-SL} \textit{versus} aleatório e\\
         envelope (Seção~\ref{sec:res-l0-drisl}), reexecução\\
         independente (Seção~\ref{sec:res-l0-replay})};

    \draw[seta] (aleat.east) -- (proto.150);
    \draw[seta] (drisl.east) -- (proto.west);
    \draw[seta] (ag.east) -- (proto.210);
    \draw[seta] (proto) -- (leit);

    % ------------- anti-circularidade do envelope --------------------------
    \node[rot, text width=44mm, anchor=north] (anticirc) at (0.6,-7.5)
        {anti-circularidade do envelope: a aptidão do AG usa uma partição de
         aferição disjunta do teste; os indivíduos finais são reavaliados no
         conjunto de teste intocado, e é esse valor, tipicamente inferior,
         que os resultados reportam};
    \draw[seta, dashed] (anticirc.north) -- (ag.south);
\end{tikzpicture}

% Uso standalone (pré-visualização):
%   \documentclass[tikz,border=6pt]{standalone}
%   \usetikzlibrary{arrows.meta, positioning}
%   \begin{document}% ============================================================================
% ESQUEMA PROPOSTO (banca, sugestão — não mergeado): o desenho da avaliação
% da composição do L0 e da partida a frio (P1/P2 no vocabulário da defesa):
% três origens para o conjunto inicial, um único protocolo de medição, e a
% nota de anti-circularidade do envelope. Ilustra a Seção sec:metodo-l0 e
% complementa o esq-drisl (que mostra o algoritmo; este mostra a avaliação).
%
% Uso no Cap. 3 (dentro de um ambiente figure):
%   \input{3-metodo/esquemas-propostos/esq-l0-tres-origens}
% Uso standalone (pré-visualização):
%   \documentclass[tikz,border=6pt]{standalone}
%   \usetikzlibrary{arrows.meta, positioning}
%   \begin{document}\input{esq-l0-tres-origens.tex}\end{document}
%   (no modo standalone, troque as \ref{...} por texto fixo ou defina-as)
% Requer apenas arrows.meta e positioning (já em packages.tex). Legível em
% P&B: origem = branco; protocolo = cinza claro; leitura = cinza médio.
% FREEZE respeitado: todos os números são os da própria Seção 3.6
% (47 tamanhos, 10 a 200.000, 30 repetições; 4 cenários, 10 tamanhos,
% 10 a 30.000), nenhum resultado afirmado — a leitura remete ao Cap. 4.
% ============================================================================
\begin{tikzpicture}[
    x=1cm, y=1cm,
    caixa/.style={draw, rounded corners=2pt, align=center, text width=42mm,
                  minimum height=11mm, inner sep=3pt, font=\footnotesize},
    protocolo/.style={caixa, fill=gray!8},
    leitura/.style={caixa, fill=gray!14},
    seta/.style={-{Stealth[length=2.4mm]}, thick},
    rot/.style={font=\scriptsize, align=center, inner sep=2pt}
]
    % ------------- três origens para o conjunto inicial --------------------
    % (o AG fica na linha de baixo, vizinho da sua nota de anti-circularidade:
    % com ele no meio, a seta tracejada atravessava a caixa do DRI-SL)
    \node[caixa, text width=44mm] (aleat) at (0.6,0)
        {\textbf{amostragem aleatória}\\(dispersão do sorteio):\\
         47 tamanhos de $L_0$, de 10 a 200.000; 30 repetições por tamanho};
    \node[caixa, text width=44mm] (drisl) at (0.6,-2.9)
        {\textbf{DRI-SL} (proposto):\\determinístico e sem acesso a
         rótulos; algoritmo na Seção~\ref{sec:metodo-drisl}};
    \node[caixa, text width=44mm] (ag) at (0.6,-5.8)
        {\textbf{algoritmo genético}\\(envelope: teto e piso):\\
         4 cenários (max/min de Acurácia\\e de Macro F1) em 10 tamanhos,\\
         de 10 a 30.000};

    % ------------- protocolo único e leitura -------------------------------
    \node[protocolo, text width=46mm] (proto) at (6.8,-2.9)
        {protocolo comum a toda origem:\\o classificador PVBin é treinado\\
         no $L_0$; Acurácia e Macro F1\\medidas no conjunto de teste};
    \node[leitura, text width=46mm] (leit) at (6.8,-6.4)
        {leitura nos resultados:\\
         sensibilidade (Seção~\ref{sec:res-l0-sens}),\\
         envelope (Seção~\ref{sec:res-l0-ag}),\\
         \mbox{DRI-SL} \textit{versus} aleatório e\\
         envelope (Seção~\ref{sec:res-l0-drisl}), reexecução\\
         independente (Seção~\ref{sec:res-l0-replay})};

    \draw[seta] (aleat.east) -- (proto.150);
    \draw[seta] (drisl.east) -- (proto.west);
    \draw[seta] (ag.east) -- (proto.210);
    \draw[seta] (proto) -- (leit);

    % ------------- anti-circularidade do envelope --------------------------
    \node[rot, text width=44mm, anchor=north] (anticirc) at (0.6,-7.5)
        {anti-circularidade do envelope: a aptidão do AG usa uma partição de
         aferição disjunta do teste; os indivíduos finais são reavaliados no
         conjunto de teste intocado, e é esse valor, tipicamente inferior,
         que os resultados reportam};
    \draw[seta, dashed] (anticirc.north) -- (ag.south);
\end{tikzpicture}
\end{document}
%   (no modo standalone, troque as \ref{...} por texto fixo ou defina-as)
% Requer apenas arrows.meta e positioning (já em packages.tex). Legível em
% P&B: origem = branco; protocolo = cinza claro; leitura = cinza médio.
% FREEZE respeitado: todos os números são os da própria Seção 3.6
% (47 tamanhos, 10 a 200.000, 30 repetições; 4 cenários, 10 tamanhos,
% 10 a 30.000), nenhum resultado afirmado — a leitura remete ao Cap. 4.
% ============================================================================
\begin{tikzpicture}[
    x=1cm, y=1cm,
    caixa/.style={draw, rounded corners=2pt, align=center, text width=42mm,
                  minimum height=11mm, inner sep=3pt, font=\footnotesize},
    protocolo/.style={caixa, fill=gray!8},
    leitura/.style={caixa, fill=gray!14},
    seta/.style={-{Stealth[length=2.4mm]}, thick},
    rot/.style={font=\scriptsize, align=center, inner sep=2pt}
]
    % ------------- três origens para o conjunto inicial --------------------
    % (o AG fica na linha de baixo, vizinho da sua nota de anti-circularidade:
    % com ele no meio, a seta tracejada atravessava a caixa do DRI-SL)
    \node[caixa, text width=44mm] (aleat) at (0.6,0)
        {\textbf{amostragem aleatória}\\(dispersão do sorteio):\\
         47 tamanhos de $L_0$, de 10 a 200.000; 30 repetições por tamanho};
    \node[caixa, text width=44mm] (drisl) at (0.6,-2.9)
        {\textbf{DRI-SL} (proposto):\\determinístico e sem acesso a
         rótulos; algoritmo na Seção~\ref{sec:metodo-drisl}};
    \node[caixa, text width=44mm] (ag) at (0.6,-5.8)
        {\textbf{algoritmo genético}\\(envelope: teto e piso):\\
         4 cenários (max/min de Acurácia\\e de Macro F1) em 10 tamanhos,\\
         de 10 a 30.000};

    % ------------- protocolo único e leitura -------------------------------
    \node[protocolo, text width=46mm] (proto) at (6.8,-2.9)
        {protocolo comum a toda origem:\\o classificador PVBin é treinado\\
         no $L_0$; Acurácia e Macro F1\\medidas no conjunto de teste};
    \node[leitura, text width=46mm] (leit) at (6.8,-6.4)
        {leitura nos resultados:\\
         sensibilidade (Seção~\ref{sec:res-l0-sens}),\\
         envelope (Seção~\ref{sec:res-l0-ag}),\\
         \mbox{DRI-SL} \textit{versus} aleatório e\\
         envelope (Seção~\ref{sec:res-l0-drisl}), reexecução\\
         independente (Seção~\ref{sec:res-l0-replay})};

    \draw[seta] (aleat.east) -- (proto.150);
    \draw[seta] (drisl.east) -- (proto.west);
    \draw[seta] (ag.east) -- (proto.210);
    \draw[seta] (proto) -- (leit);

    % ------------- anti-circularidade do envelope --------------------------
    \node[rot, text width=44mm, anchor=north] (anticirc) at (0.6,-7.5)
        {anti-circularidade do envelope: a aptidão do AG usa uma partição de
         aferição disjunta do teste; os indivíduos finais são reavaliados no
         conjunto de teste intocado, e é esse valor, tipicamente inferior,
         que os resultados reportam};
    \draw[seta, dashed] (anticirc.north) -- (ag.south);
\end{tikzpicture}

% Uso standalone (pré-visualização):
%   \documentclass[tikz,border=6pt]{standalone}
%   \usetikzlibrary{arrows.meta, positioning}
%   \begin{document}% ============================================================================
% ESQUEMA PROPOSTO (banca, sugestão — não mergeado): o desenho da avaliação
% da composição do L0 e da partida a frio (P1/P2 no vocabulário da defesa):
% três origens para o conjunto inicial, um único protocolo de medição, e a
% nota de anti-circularidade do envelope. Ilustra a Seção sec:metodo-l0 e
% complementa o esq-drisl (que mostra o algoritmo; este mostra a avaliação).
%
% Uso no Cap. 3 (dentro de um ambiente figure):
%   % ============================================================================
% ESQUEMA PROPOSTO (banca, sugestão — não mergeado): o desenho da avaliação
% da composição do L0 e da partida a frio (P1/P2 no vocabulário da defesa):
% três origens para o conjunto inicial, um único protocolo de medição, e a
% nota de anti-circularidade do envelope. Ilustra a Seção sec:metodo-l0 e
% complementa o esq-drisl (que mostra o algoritmo; este mostra a avaliação).
%
% Uso no Cap. 3 (dentro de um ambiente figure):
%   \input{3-metodo/esquemas-propostos/esq-l0-tres-origens}
% Uso standalone (pré-visualização):
%   \documentclass[tikz,border=6pt]{standalone}
%   \usetikzlibrary{arrows.meta, positioning}
%   \begin{document}\input{esq-l0-tres-origens.tex}\end{document}
%   (no modo standalone, troque as \ref{...} por texto fixo ou defina-as)
% Requer apenas arrows.meta e positioning (já em packages.tex). Legível em
% P&B: origem = branco; protocolo = cinza claro; leitura = cinza médio.
% FREEZE respeitado: todos os números são os da própria Seção 3.6
% (47 tamanhos, 10 a 200.000, 30 repetições; 4 cenários, 10 tamanhos,
% 10 a 30.000), nenhum resultado afirmado — a leitura remete ao Cap. 4.
% ============================================================================
\begin{tikzpicture}[
    x=1cm, y=1cm,
    caixa/.style={draw, rounded corners=2pt, align=center, text width=42mm,
                  minimum height=11mm, inner sep=3pt, font=\footnotesize},
    protocolo/.style={caixa, fill=gray!8},
    leitura/.style={caixa, fill=gray!14},
    seta/.style={-{Stealth[length=2.4mm]}, thick},
    rot/.style={font=\scriptsize, align=center, inner sep=2pt}
]
    % ------------- três origens para o conjunto inicial --------------------
    % (o AG fica na linha de baixo, vizinho da sua nota de anti-circularidade:
    % com ele no meio, a seta tracejada atravessava a caixa do DRI-SL)
    \node[caixa, text width=44mm] (aleat) at (0.6,0)
        {\textbf{amostragem aleatória}\\(dispersão do sorteio):\\
         47 tamanhos de $L_0$, de 10 a 200.000; 30 repetições por tamanho};
    \node[caixa, text width=44mm] (drisl) at (0.6,-2.9)
        {\textbf{DRI-SL} (proposto):\\determinístico e sem acesso a
         rótulos; algoritmo na Seção~\ref{sec:metodo-drisl}};
    \node[caixa, text width=44mm] (ag) at (0.6,-5.8)
        {\textbf{algoritmo genético}\\(envelope: teto e piso):\\
         4 cenários (max/min de Acurácia\\e de Macro F1) em 10 tamanhos,\\
         de 10 a 30.000};

    % ------------- protocolo único e leitura -------------------------------
    \node[protocolo, text width=46mm] (proto) at (6.8,-2.9)
        {protocolo comum a toda origem:\\o classificador PVBin é treinado\\
         no $L_0$; Acurácia e Macro F1\\medidas no conjunto de teste};
    \node[leitura, text width=46mm] (leit) at (6.8,-6.4)
        {leitura nos resultados:\\
         sensibilidade (Seção~\ref{sec:res-l0-sens}),\\
         envelope (Seção~\ref{sec:res-l0-ag}),\\
         \mbox{DRI-SL} \textit{versus} aleatório e\\
         envelope (Seção~\ref{sec:res-l0-drisl}), reexecução\\
         independente (Seção~\ref{sec:res-l0-replay})};

    \draw[seta] (aleat.east) -- (proto.150);
    \draw[seta] (drisl.east) -- (proto.west);
    \draw[seta] (ag.east) -- (proto.210);
    \draw[seta] (proto) -- (leit);

    % ------------- anti-circularidade do envelope --------------------------
    \node[rot, text width=44mm, anchor=north] (anticirc) at (0.6,-7.5)
        {anti-circularidade do envelope: a aptidão do AG usa uma partição de
         aferição disjunta do teste; os indivíduos finais são reavaliados no
         conjunto de teste intocado, e é esse valor, tipicamente inferior,
         que os resultados reportam};
    \draw[seta, dashed] (anticirc.north) -- (ag.south);
\end{tikzpicture}

% Uso standalone (pré-visualização):
%   \documentclass[tikz,border=6pt]{standalone}
%   \usetikzlibrary{arrows.meta, positioning}
%   \begin{document}% ============================================================================
% ESQUEMA PROPOSTO (banca, sugestão — não mergeado): o desenho da avaliação
% da composição do L0 e da partida a frio (P1/P2 no vocabulário da defesa):
% três origens para o conjunto inicial, um único protocolo de medição, e a
% nota de anti-circularidade do envelope. Ilustra a Seção sec:metodo-l0 e
% complementa o esq-drisl (que mostra o algoritmo; este mostra a avaliação).
%
% Uso no Cap. 3 (dentro de um ambiente figure):
%   \input{3-metodo/esquemas-propostos/esq-l0-tres-origens}
% Uso standalone (pré-visualização):
%   \documentclass[tikz,border=6pt]{standalone}
%   \usetikzlibrary{arrows.meta, positioning}
%   \begin{document}\input{esq-l0-tres-origens.tex}\end{document}
%   (no modo standalone, troque as \ref{...} por texto fixo ou defina-as)
% Requer apenas arrows.meta e positioning (já em packages.tex). Legível em
% P&B: origem = branco; protocolo = cinza claro; leitura = cinza médio.
% FREEZE respeitado: todos os números são os da própria Seção 3.6
% (47 tamanhos, 10 a 200.000, 30 repetições; 4 cenários, 10 tamanhos,
% 10 a 30.000), nenhum resultado afirmado — a leitura remete ao Cap. 4.
% ============================================================================
\begin{tikzpicture}[
    x=1cm, y=1cm,
    caixa/.style={draw, rounded corners=2pt, align=center, text width=42mm,
                  minimum height=11mm, inner sep=3pt, font=\footnotesize},
    protocolo/.style={caixa, fill=gray!8},
    leitura/.style={caixa, fill=gray!14},
    seta/.style={-{Stealth[length=2.4mm]}, thick},
    rot/.style={font=\scriptsize, align=center, inner sep=2pt}
]
    % ------------- três origens para o conjunto inicial --------------------
    % (o AG fica na linha de baixo, vizinho da sua nota de anti-circularidade:
    % com ele no meio, a seta tracejada atravessava a caixa do DRI-SL)
    \node[caixa, text width=44mm] (aleat) at (0.6,0)
        {\textbf{amostragem aleatória}\\(dispersão do sorteio):\\
         47 tamanhos de $L_0$, de 10 a 200.000; 30 repetições por tamanho};
    \node[caixa, text width=44mm] (drisl) at (0.6,-2.9)
        {\textbf{DRI-SL} (proposto):\\determinístico e sem acesso a
         rótulos; algoritmo na Seção~\ref{sec:metodo-drisl}};
    \node[caixa, text width=44mm] (ag) at (0.6,-5.8)
        {\textbf{algoritmo genético}\\(envelope: teto e piso):\\
         4 cenários (max/min de Acurácia\\e de Macro F1) em 10 tamanhos,\\
         de 10 a 30.000};

    % ------------- protocolo único e leitura -------------------------------
    \node[protocolo, text width=46mm] (proto) at (6.8,-2.9)
        {protocolo comum a toda origem:\\o classificador PVBin é treinado\\
         no $L_0$; Acurácia e Macro F1\\medidas no conjunto de teste};
    \node[leitura, text width=46mm] (leit) at (6.8,-6.4)
        {leitura nos resultados:\\
         sensibilidade (Seção~\ref{sec:res-l0-sens}),\\
         envelope (Seção~\ref{sec:res-l0-ag}),\\
         \mbox{DRI-SL} \textit{versus} aleatório e\\
         envelope (Seção~\ref{sec:res-l0-drisl}), reexecução\\
         independente (Seção~\ref{sec:res-l0-replay})};

    \draw[seta] (aleat.east) -- (proto.150);
    \draw[seta] (drisl.east) -- (proto.west);
    \draw[seta] (ag.east) -- (proto.210);
    \draw[seta] (proto) -- (leit);

    % ------------- anti-circularidade do envelope --------------------------
    \node[rot, text width=44mm, anchor=north] (anticirc) at (0.6,-7.5)
        {anti-circularidade do envelope: a aptidão do AG usa uma partição de
         aferição disjunta do teste; os indivíduos finais são reavaliados no
         conjunto de teste intocado, e é esse valor, tipicamente inferior,
         que os resultados reportam};
    \draw[seta, dashed] (anticirc.north) -- (ag.south);
\end{tikzpicture}
\end{document}
%   (no modo standalone, troque as \ref{...} por texto fixo ou defina-as)
% Requer apenas arrows.meta e positioning (já em packages.tex). Legível em
% P&B: origem = branco; protocolo = cinza claro; leitura = cinza médio.
% FREEZE respeitado: todos os números são os da própria Seção 3.6
% (47 tamanhos, 10 a 200.000, 30 repetições; 4 cenários, 10 tamanhos,
% 10 a 30.000), nenhum resultado afirmado — a leitura remete ao Cap. 4.
% ============================================================================
\begin{tikzpicture}[
    x=1cm, y=1cm,
    caixa/.style={draw, rounded corners=2pt, align=center, text width=42mm,
                  minimum height=11mm, inner sep=3pt, font=\footnotesize},
    protocolo/.style={caixa, fill=gray!8},
    leitura/.style={caixa, fill=gray!14},
    seta/.style={-{Stealth[length=2.4mm]}, thick},
    rot/.style={font=\scriptsize, align=center, inner sep=2pt}
]
    % ------------- três origens para o conjunto inicial --------------------
    % (o AG fica na linha de baixo, vizinho da sua nota de anti-circularidade:
    % com ele no meio, a seta tracejada atravessava a caixa do DRI-SL)
    \node[caixa, text width=44mm] (aleat) at (0.6,0)
        {\textbf{amostragem aleatória}\\(dispersão do sorteio):\\
         47 tamanhos de $L_0$, de 10 a 200.000; 30 repetições por tamanho};
    \node[caixa, text width=44mm] (drisl) at (0.6,-2.9)
        {\textbf{DRI-SL} (proposto):\\determinístico e sem acesso a
         rótulos; algoritmo na Seção~\ref{sec:metodo-drisl}};
    \node[caixa, text width=44mm] (ag) at (0.6,-5.8)
        {\textbf{algoritmo genético}\\(envelope: teto e piso):\\
         4 cenários (max/min de Acurácia\\e de Macro F1) em 10 tamanhos,\\
         de 10 a 30.000};

    % ------------- protocolo único e leitura -------------------------------
    \node[protocolo, text width=46mm] (proto) at (6.8,-2.9)
        {protocolo comum a toda origem:\\o classificador PVBin é treinado\\
         no $L_0$; Acurácia e Macro F1\\medidas no conjunto de teste};
    \node[leitura, text width=46mm] (leit) at (6.8,-6.4)
        {leitura nos resultados:\\
         sensibilidade (Seção~\ref{sec:res-l0-sens}),\\
         envelope (Seção~\ref{sec:res-l0-ag}),\\
         \mbox{DRI-SL} \textit{versus} aleatório e\\
         envelope (Seção~\ref{sec:res-l0-drisl}), reexecução\\
         independente (Seção~\ref{sec:res-l0-replay})};

    \draw[seta] (aleat.east) -- (proto.150);
    \draw[seta] (drisl.east) -- (proto.west);
    \draw[seta] (ag.east) -- (proto.210);
    \draw[seta] (proto) -- (leit);

    % ------------- anti-circularidade do envelope --------------------------
    \node[rot, text width=44mm, anchor=north] (anticirc) at (0.6,-7.5)
        {anti-circularidade do envelope: a aptidão do AG usa uma partição de
         aferição disjunta do teste; os indivíduos finais são reavaliados no
         conjunto de teste intocado, e é esse valor, tipicamente inferior,
         que os resultados reportam};
    \draw[seta, dashed] (anticirc.north) -- (ag.south);
\end{tikzpicture}
\end{document}
%   (no modo standalone, troque as \ref{...} por texto fixo ou defina-as)
% Requer apenas arrows.meta e positioning (já em packages.tex). Legível em
% P&B: origem = branco; protocolo = cinza claro; leitura = cinza médio.
% FREEZE respeitado: todos os números são os da própria Seção 3.6
% (47 tamanhos, 10 a 200.000, 30 repetições; 4 cenários, 10 tamanhos,
% 10 a 30.000), nenhum resultado afirmado — a leitura remete ao Cap. 4.
% ============================================================================
\begin{tikzpicture}[
    x=1cm, y=1cm,
    caixa/.style={draw, rounded corners=2pt, align=center, text width=42mm,
                  minimum height=11mm, inner sep=3pt, font=\footnotesize},
    protocolo/.style={caixa, fill=gray!8},
    leitura/.style={caixa, fill=gray!14},
    seta/.style={-{Stealth[length=2.4mm]}, thick},
    rot/.style={font=\scriptsize, align=center, inner sep=2pt}
]
    % ------------- três origens para o conjunto inicial --------------------
    % (o AG fica na linha de baixo, vizinho da sua nota de anti-circularidade:
    % com ele no meio, a seta tracejada atravessava a caixa do DRI-SL)
    \node[caixa, text width=44mm] (aleat) at (0.6,0)
        {\textbf{amostragem aleatória}\\(dispersão do sorteio):\\
         47 tamanhos de $L_0$, de 10 a 200.000; 30 repetições por tamanho};
    \node[caixa, text width=44mm] (drisl) at (0.6,-2.9)
        {\textbf{DRI-SL} (proposto):\\determinístico e sem acesso a
         rótulos; algoritmo na Seção~\ref{sec:metodo-drisl}};
    \node[caixa, text width=44mm] (ag) at (0.6,-5.8)
        {\textbf{algoritmo genético}\\(envelope: teto e piso):\\
         4 cenários (max/min de Acurácia\\e de Macro F1) em 10 tamanhos,\\
         de 10 a 30.000};

    % ------------- protocolo único e leitura -------------------------------
    \node[protocolo, text width=46mm] (proto) at (6.8,-2.9)
        {protocolo comum a toda origem:\\o classificador PVBin é treinado\\
         no $L_0$; Acurácia e Macro F1\\medidas no conjunto de teste};
    \node[leitura, text width=46mm] (leit) at (6.8,-6.4)
        {leitura nos resultados:\\
         sensibilidade (Seção~\ref{sec:res-l0-sens}),\\
         envelope (Seção~\ref{sec:res-l0-ag}),\\
         \mbox{DRI-SL} \textit{versus} aleatório e\\
         envelope (Seção~\ref{sec:res-l0-drisl}), reexecução\\
         independente (Seção~\ref{sec:res-l0-replay})};

    \draw[seta] (aleat.east) -- (proto.150);
    \draw[seta] (drisl.east) -- (proto.west);
    \draw[seta] (ag.east) -- (proto.210);
    \draw[seta] (proto) -- (leit);

    % ------------- anti-circularidade do envelope --------------------------
    \node[rot, text width=44mm, anchor=north] (anticirc) at (0.6,-7.5)
        {anti-circularidade do envelope: a aptidão do AG usa uma partição de
         aferição disjunta do teste; os indivíduos finais são reavaliados no
         conjunto de teste intocado, e é esse valor, tipicamente inferior,
         que os resultados reportam};
    \draw[seta, dashed] (anticirc.north) -- (ag.south);
\end{tikzpicture}

% Uso standalone (pré-visualização):
%   \documentclass[tikz,border=6pt]{standalone}
%   \usetikzlibrary{arrows.meta, positioning}
%   \begin{document}% ============================================================================
% ESQUEMA PROPOSTO (banca, sugestão — não mergeado): o desenho da avaliação
% da composição do L0 e da partida a frio (P1/P2 no vocabulário da defesa):
% três origens para o conjunto inicial, um único protocolo de medição, e a
% nota de anti-circularidade do envelope. Ilustra a Seção sec:metodo-l0 e
% complementa o esq-drisl (que mostra o algoritmo; este mostra a avaliação).
%
% Uso no Cap. 3 (dentro de um ambiente figure):
%   % ============================================================================
% ESQUEMA PROPOSTO (banca, sugestão — não mergeado): o desenho da avaliação
% da composição do L0 e da partida a frio (P1/P2 no vocabulário da defesa):
% três origens para o conjunto inicial, um único protocolo de medição, e a
% nota de anti-circularidade do envelope. Ilustra a Seção sec:metodo-l0 e
% complementa o esq-drisl (que mostra o algoritmo; este mostra a avaliação).
%
% Uso no Cap. 3 (dentro de um ambiente figure):
%   % ============================================================================
% ESQUEMA PROPOSTO (banca, sugestão — não mergeado): o desenho da avaliação
% da composição do L0 e da partida a frio (P1/P2 no vocabulário da defesa):
% três origens para o conjunto inicial, um único protocolo de medição, e a
% nota de anti-circularidade do envelope. Ilustra a Seção sec:metodo-l0 e
% complementa o esq-drisl (que mostra o algoritmo; este mostra a avaliação).
%
% Uso no Cap. 3 (dentro de um ambiente figure):
%   \input{3-metodo/esquemas-propostos/esq-l0-tres-origens}
% Uso standalone (pré-visualização):
%   \documentclass[tikz,border=6pt]{standalone}
%   \usetikzlibrary{arrows.meta, positioning}
%   \begin{document}\input{esq-l0-tres-origens.tex}\end{document}
%   (no modo standalone, troque as \ref{...} por texto fixo ou defina-as)
% Requer apenas arrows.meta e positioning (já em packages.tex). Legível em
% P&B: origem = branco; protocolo = cinza claro; leitura = cinza médio.
% FREEZE respeitado: todos os números são os da própria Seção 3.6
% (47 tamanhos, 10 a 200.000, 30 repetições; 4 cenários, 10 tamanhos,
% 10 a 30.000), nenhum resultado afirmado — a leitura remete ao Cap. 4.
% ============================================================================
\begin{tikzpicture}[
    x=1cm, y=1cm,
    caixa/.style={draw, rounded corners=2pt, align=center, text width=42mm,
                  minimum height=11mm, inner sep=3pt, font=\footnotesize},
    protocolo/.style={caixa, fill=gray!8},
    leitura/.style={caixa, fill=gray!14},
    seta/.style={-{Stealth[length=2.4mm]}, thick},
    rot/.style={font=\scriptsize, align=center, inner sep=2pt}
]
    % ------------- três origens para o conjunto inicial --------------------
    % (o AG fica na linha de baixo, vizinho da sua nota de anti-circularidade:
    % com ele no meio, a seta tracejada atravessava a caixa do DRI-SL)
    \node[caixa, text width=44mm] (aleat) at (0.6,0)
        {\textbf{amostragem aleatória}\\(dispersão do sorteio):\\
         47 tamanhos de $L_0$, de 10 a 200.000; 30 repetições por tamanho};
    \node[caixa, text width=44mm] (drisl) at (0.6,-2.9)
        {\textbf{DRI-SL} (proposto):\\determinístico e sem acesso a
         rótulos; algoritmo na Seção~\ref{sec:metodo-drisl}};
    \node[caixa, text width=44mm] (ag) at (0.6,-5.8)
        {\textbf{algoritmo genético}\\(envelope: teto e piso):\\
         4 cenários (max/min de Acurácia\\e de Macro F1) em 10 tamanhos,\\
         de 10 a 30.000};

    % ------------- protocolo único e leitura -------------------------------
    \node[protocolo, text width=46mm] (proto) at (6.8,-2.9)
        {protocolo comum a toda origem:\\o classificador PVBin é treinado\\
         no $L_0$; Acurácia e Macro F1\\medidas no conjunto de teste};
    \node[leitura, text width=46mm] (leit) at (6.8,-6.4)
        {leitura nos resultados:\\
         sensibilidade (Seção~\ref{sec:res-l0-sens}),\\
         envelope (Seção~\ref{sec:res-l0-ag}),\\
         \mbox{DRI-SL} \textit{versus} aleatório e\\
         envelope (Seção~\ref{sec:res-l0-drisl}), reexecução\\
         independente (Seção~\ref{sec:res-l0-replay})};

    \draw[seta] (aleat.east) -- (proto.150);
    \draw[seta] (drisl.east) -- (proto.west);
    \draw[seta] (ag.east) -- (proto.210);
    \draw[seta] (proto) -- (leit);

    % ------------- anti-circularidade do envelope --------------------------
    \node[rot, text width=44mm, anchor=north] (anticirc) at (0.6,-7.5)
        {anti-circularidade do envelope: a aptidão do AG usa uma partição de
         aferição disjunta do teste; os indivíduos finais são reavaliados no
         conjunto de teste intocado, e é esse valor, tipicamente inferior,
         que os resultados reportam};
    \draw[seta, dashed] (anticirc.north) -- (ag.south);
\end{tikzpicture}

% Uso standalone (pré-visualização):
%   \documentclass[tikz,border=6pt]{standalone}
%   \usetikzlibrary{arrows.meta, positioning}
%   \begin{document}% ============================================================================
% ESQUEMA PROPOSTO (banca, sugestão — não mergeado): o desenho da avaliação
% da composição do L0 e da partida a frio (P1/P2 no vocabulário da defesa):
% três origens para o conjunto inicial, um único protocolo de medição, e a
% nota de anti-circularidade do envelope. Ilustra a Seção sec:metodo-l0 e
% complementa o esq-drisl (que mostra o algoritmo; este mostra a avaliação).
%
% Uso no Cap. 3 (dentro de um ambiente figure):
%   \input{3-metodo/esquemas-propostos/esq-l0-tres-origens}
% Uso standalone (pré-visualização):
%   \documentclass[tikz,border=6pt]{standalone}
%   \usetikzlibrary{arrows.meta, positioning}
%   \begin{document}\input{esq-l0-tres-origens.tex}\end{document}
%   (no modo standalone, troque as \ref{...} por texto fixo ou defina-as)
% Requer apenas arrows.meta e positioning (já em packages.tex). Legível em
% P&B: origem = branco; protocolo = cinza claro; leitura = cinza médio.
% FREEZE respeitado: todos os números são os da própria Seção 3.6
% (47 tamanhos, 10 a 200.000, 30 repetições; 4 cenários, 10 tamanhos,
% 10 a 30.000), nenhum resultado afirmado — a leitura remete ao Cap. 4.
% ============================================================================
\begin{tikzpicture}[
    x=1cm, y=1cm,
    caixa/.style={draw, rounded corners=2pt, align=center, text width=42mm,
                  minimum height=11mm, inner sep=3pt, font=\footnotesize},
    protocolo/.style={caixa, fill=gray!8},
    leitura/.style={caixa, fill=gray!14},
    seta/.style={-{Stealth[length=2.4mm]}, thick},
    rot/.style={font=\scriptsize, align=center, inner sep=2pt}
]
    % ------------- três origens para o conjunto inicial --------------------
    % (o AG fica na linha de baixo, vizinho da sua nota de anti-circularidade:
    % com ele no meio, a seta tracejada atravessava a caixa do DRI-SL)
    \node[caixa, text width=44mm] (aleat) at (0.6,0)
        {\textbf{amostragem aleatória}\\(dispersão do sorteio):\\
         47 tamanhos de $L_0$, de 10 a 200.000; 30 repetições por tamanho};
    \node[caixa, text width=44mm] (drisl) at (0.6,-2.9)
        {\textbf{DRI-SL} (proposto):\\determinístico e sem acesso a
         rótulos; algoritmo na Seção~\ref{sec:metodo-drisl}};
    \node[caixa, text width=44mm] (ag) at (0.6,-5.8)
        {\textbf{algoritmo genético}\\(envelope: teto e piso):\\
         4 cenários (max/min de Acurácia\\e de Macro F1) em 10 tamanhos,\\
         de 10 a 30.000};

    % ------------- protocolo único e leitura -------------------------------
    \node[protocolo, text width=46mm] (proto) at (6.8,-2.9)
        {protocolo comum a toda origem:\\o classificador PVBin é treinado\\
         no $L_0$; Acurácia e Macro F1\\medidas no conjunto de teste};
    \node[leitura, text width=46mm] (leit) at (6.8,-6.4)
        {leitura nos resultados:\\
         sensibilidade (Seção~\ref{sec:res-l0-sens}),\\
         envelope (Seção~\ref{sec:res-l0-ag}),\\
         \mbox{DRI-SL} \textit{versus} aleatório e\\
         envelope (Seção~\ref{sec:res-l0-drisl}), reexecução\\
         independente (Seção~\ref{sec:res-l0-replay})};

    \draw[seta] (aleat.east) -- (proto.150);
    \draw[seta] (drisl.east) -- (proto.west);
    \draw[seta] (ag.east) -- (proto.210);
    \draw[seta] (proto) -- (leit);

    % ------------- anti-circularidade do envelope --------------------------
    \node[rot, text width=44mm, anchor=north] (anticirc) at (0.6,-7.5)
        {anti-circularidade do envelope: a aptidão do AG usa uma partição de
         aferição disjunta do teste; os indivíduos finais são reavaliados no
         conjunto de teste intocado, e é esse valor, tipicamente inferior,
         que os resultados reportam};
    \draw[seta, dashed] (anticirc.north) -- (ag.south);
\end{tikzpicture}
\end{document}
%   (no modo standalone, troque as \ref{...} por texto fixo ou defina-as)
% Requer apenas arrows.meta e positioning (já em packages.tex). Legível em
% P&B: origem = branco; protocolo = cinza claro; leitura = cinza médio.
% FREEZE respeitado: todos os números são os da própria Seção 3.6
% (47 tamanhos, 10 a 200.000, 30 repetições; 4 cenários, 10 tamanhos,
% 10 a 30.000), nenhum resultado afirmado — a leitura remete ao Cap. 4.
% ============================================================================
\begin{tikzpicture}[
    x=1cm, y=1cm,
    caixa/.style={draw, rounded corners=2pt, align=center, text width=42mm,
                  minimum height=11mm, inner sep=3pt, font=\footnotesize},
    protocolo/.style={caixa, fill=gray!8},
    leitura/.style={caixa, fill=gray!14},
    seta/.style={-{Stealth[length=2.4mm]}, thick},
    rot/.style={font=\scriptsize, align=center, inner sep=2pt}
]
    % ------------- três origens para o conjunto inicial --------------------
    % (o AG fica na linha de baixo, vizinho da sua nota de anti-circularidade:
    % com ele no meio, a seta tracejada atravessava a caixa do DRI-SL)
    \node[caixa, text width=44mm] (aleat) at (0.6,0)
        {\textbf{amostragem aleatória}\\(dispersão do sorteio):\\
         47 tamanhos de $L_0$, de 10 a 200.000; 30 repetições por tamanho};
    \node[caixa, text width=44mm] (drisl) at (0.6,-2.9)
        {\textbf{DRI-SL} (proposto):\\determinístico e sem acesso a
         rótulos; algoritmo na Seção~\ref{sec:metodo-drisl}};
    \node[caixa, text width=44mm] (ag) at (0.6,-5.8)
        {\textbf{algoritmo genético}\\(envelope: teto e piso):\\
         4 cenários (max/min de Acurácia\\e de Macro F1) em 10 tamanhos,\\
         de 10 a 30.000};

    % ------------- protocolo único e leitura -------------------------------
    \node[protocolo, text width=46mm] (proto) at (6.8,-2.9)
        {protocolo comum a toda origem:\\o classificador PVBin é treinado\\
         no $L_0$; Acurácia e Macro F1\\medidas no conjunto de teste};
    \node[leitura, text width=46mm] (leit) at (6.8,-6.4)
        {leitura nos resultados:\\
         sensibilidade (Seção~\ref{sec:res-l0-sens}),\\
         envelope (Seção~\ref{sec:res-l0-ag}),\\
         \mbox{DRI-SL} \textit{versus} aleatório e\\
         envelope (Seção~\ref{sec:res-l0-drisl}), reexecução\\
         independente (Seção~\ref{sec:res-l0-replay})};

    \draw[seta] (aleat.east) -- (proto.150);
    \draw[seta] (drisl.east) -- (proto.west);
    \draw[seta] (ag.east) -- (proto.210);
    \draw[seta] (proto) -- (leit);

    % ------------- anti-circularidade do envelope --------------------------
    \node[rot, text width=44mm, anchor=north] (anticirc) at (0.6,-7.5)
        {anti-circularidade do envelope: a aptidão do AG usa uma partição de
         aferição disjunta do teste; os indivíduos finais são reavaliados no
         conjunto de teste intocado, e é esse valor, tipicamente inferior,
         que os resultados reportam};
    \draw[seta, dashed] (anticirc.north) -- (ag.south);
\end{tikzpicture}

% Uso standalone (pré-visualização):
%   \documentclass[tikz,border=6pt]{standalone}
%   \usetikzlibrary{arrows.meta, positioning}
%   \begin{document}% ============================================================================
% ESQUEMA PROPOSTO (banca, sugestão — não mergeado): o desenho da avaliação
% da composição do L0 e da partida a frio (P1/P2 no vocabulário da defesa):
% três origens para o conjunto inicial, um único protocolo de medição, e a
% nota de anti-circularidade do envelope. Ilustra a Seção sec:metodo-l0 e
% complementa o esq-drisl (que mostra o algoritmo; este mostra a avaliação).
%
% Uso no Cap. 3 (dentro de um ambiente figure):
%   % ============================================================================
% ESQUEMA PROPOSTO (banca, sugestão — não mergeado): o desenho da avaliação
% da composição do L0 e da partida a frio (P1/P2 no vocabulário da defesa):
% três origens para o conjunto inicial, um único protocolo de medição, e a
% nota de anti-circularidade do envelope. Ilustra a Seção sec:metodo-l0 e
% complementa o esq-drisl (que mostra o algoritmo; este mostra a avaliação).
%
% Uso no Cap. 3 (dentro de um ambiente figure):
%   \input{3-metodo/esquemas-propostos/esq-l0-tres-origens}
% Uso standalone (pré-visualização):
%   \documentclass[tikz,border=6pt]{standalone}
%   \usetikzlibrary{arrows.meta, positioning}
%   \begin{document}\input{esq-l0-tres-origens.tex}\end{document}
%   (no modo standalone, troque as \ref{...} por texto fixo ou defina-as)
% Requer apenas arrows.meta e positioning (já em packages.tex). Legível em
% P&B: origem = branco; protocolo = cinza claro; leitura = cinza médio.
% FREEZE respeitado: todos os números são os da própria Seção 3.6
% (47 tamanhos, 10 a 200.000, 30 repetições; 4 cenários, 10 tamanhos,
% 10 a 30.000), nenhum resultado afirmado — a leitura remete ao Cap. 4.
% ============================================================================
\begin{tikzpicture}[
    x=1cm, y=1cm,
    caixa/.style={draw, rounded corners=2pt, align=center, text width=42mm,
                  minimum height=11mm, inner sep=3pt, font=\footnotesize},
    protocolo/.style={caixa, fill=gray!8},
    leitura/.style={caixa, fill=gray!14},
    seta/.style={-{Stealth[length=2.4mm]}, thick},
    rot/.style={font=\scriptsize, align=center, inner sep=2pt}
]
    % ------------- três origens para o conjunto inicial --------------------
    % (o AG fica na linha de baixo, vizinho da sua nota de anti-circularidade:
    % com ele no meio, a seta tracejada atravessava a caixa do DRI-SL)
    \node[caixa, text width=44mm] (aleat) at (0.6,0)
        {\textbf{amostragem aleatória}\\(dispersão do sorteio):\\
         47 tamanhos de $L_0$, de 10 a 200.000; 30 repetições por tamanho};
    \node[caixa, text width=44mm] (drisl) at (0.6,-2.9)
        {\textbf{DRI-SL} (proposto):\\determinístico e sem acesso a
         rótulos; algoritmo na Seção~\ref{sec:metodo-drisl}};
    \node[caixa, text width=44mm] (ag) at (0.6,-5.8)
        {\textbf{algoritmo genético}\\(envelope: teto e piso):\\
         4 cenários (max/min de Acurácia\\e de Macro F1) em 10 tamanhos,\\
         de 10 a 30.000};

    % ------------- protocolo único e leitura -------------------------------
    \node[protocolo, text width=46mm] (proto) at (6.8,-2.9)
        {protocolo comum a toda origem:\\o classificador PVBin é treinado\\
         no $L_0$; Acurácia e Macro F1\\medidas no conjunto de teste};
    \node[leitura, text width=46mm] (leit) at (6.8,-6.4)
        {leitura nos resultados:\\
         sensibilidade (Seção~\ref{sec:res-l0-sens}),\\
         envelope (Seção~\ref{sec:res-l0-ag}),\\
         \mbox{DRI-SL} \textit{versus} aleatório e\\
         envelope (Seção~\ref{sec:res-l0-drisl}), reexecução\\
         independente (Seção~\ref{sec:res-l0-replay})};

    \draw[seta] (aleat.east) -- (proto.150);
    \draw[seta] (drisl.east) -- (proto.west);
    \draw[seta] (ag.east) -- (proto.210);
    \draw[seta] (proto) -- (leit);

    % ------------- anti-circularidade do envelope --------------------------
    \node[rot, text width=44mm, anchor=north] (anticirc) at (0.6,-7.5)
        {anti-circularidade do envelope: a aptidão do AG usa uma partição de
         aferição disjunta do teste; os indivíduos finais são reavaliados no
         conjunto de teste intocado, e é esse valor, tipicamente inferior,
         que os resultados reportam};
    \draw[seta, dashed] (anticirc.north) -- (ag.south);
\end{tikzpicture}

% Uso standalone (pré-visualização):
%   \documentclass[tikz,border=6pt]{standalone}
%   \usetikzlibrary{arrows.meta, positioning}
%   \begin{document}% ============================================================================
% ESQUEMA PROPOSTO (banca, sugestão — não mergeado): o desenho da avaliação
% da composição do L0 e da partida a frio (P1/P2 no vocabulário da defesa):
% três origens para o conjunto inicial, um único protocolo de medição, e a
% nota de anti-circularidade do envelope. Ilustra a Seção sec:metodo-l0 e
% complementa o esq-drisl (que mostra o algoritmo; este mostra a avaliação).
%
% Uso no Cap. 3 (dentro de um ambiente figure):
%   \input{3-metodo/esquemas-propostos/esq-l0-tres-origens}
% Uso standalone (pré-visualização):
%   \documentclass[tikz,border=6pt]{standalone}
%   \usetikzlibrary{arrows.meta, positioning}
%   \begin{document}\input{esq-l0-tres-origens.tex}\end{document}
%   (no modo standalone, troque as \ref{...} por texto fixo ou defina-as)
% Requer apenas arrows.meta e positioning (já em packages.tex). Legível em
% P&B: origem = branco; protocolo = cinza claro; leitura = cinza médio.
% FREEZE respeitado: todos os números são os da própria Seção 3.6
% (47 tamanhos, 10 a 200.000, 30 repetições; 4 cenários, 10 tamanhos,
% 10 a 30.000), nenhum resultado afirmado — a leitura remete ao Cap. 4.
% ============================================================================
\begin{tikzpicture}[
    x=1cm, y=1cm,
    caixa/.style={draw, rounded corners=2pt, align=center, text width=42mm,
                  minimum height=11mm, inner sep=3pt, font=\footnotesize},
    protocolo/.style={caixa, fill=gray!8},
    leitura/.style={caixa, fill=gray!14},
    seta/.style={-{Stealth[length=2.4mm]}, thick},
    rot/.style={font=\scriptsize, align=center, inner sep=2pt}
]
    % ------------- três origens para o conjunto inicial --------------------
    % (o AG fica na linha de baixo, vizinho da sua nota de anti-circularidade:
    % com ele no meio, a seta tracejada atravessava a caixa do DRI-SL)
    \node[caixa, text width=44mm] (aleat) at (0.6,0)
        {\textbf{amostragem aleatória}\\(dispersão do sorteio):\\
         47 tamanhos de $L_0$, de 10 a 200.000; 30 repetições por tamanho};
    \node[caixa, text width=44mm] (drisl) at (0.6,-2.9)
        {\textbf{DRI-SL} (proposto):\\determinístico e sem acesso a
         rótulos; algoritmo na Seção~\ref{sec:metodo-drisl}};
    \node[caixa, text width=44mm] (ag) at (0.6,-5.8)
        {\textbf{algoritmo genético}\\(envelope: teto e piso):\\
         4 cenários (max/min de Acurácia\\e de Macro F1) em 10 tamanhos,\\
         de 10 a 30.000};

    % ------------- protocolo único e leitura -------------------------------
    \node[protocolo, text width=46mm] (proto) at (6.8,-2.9)
        {protocolo comum a toda origem:\\o classificador PVBin é treinado\\
         no $L_0$; Acurácia e Macro F1\\medidas no conjunto de teste};
    \node[leitura, text width=46mm] (leit) at (6.8,-6.4)
        {leitura nos resultados:\\
         sensibilidade (Seção~\ref{sec:res-l0-sens}),\\
         envelope (Seção~\ref{sec:res-l0-ag}),\\
         \mbox{DRI-SL} \textit{versus} aleatório e\\
         envelope (Seção~\ref{sec:res-l0-drisl}), reexecução\\
         independente (Seção~\ref{sec:res-l0-replay})};

    \draw[seta] (aleat.east) -- (proto.150);
    \draw[seta] (drisl.east) -- (proto.west);
    \draw[seta] (ag.east) -- (proto.210);
    \draw[seta] (proto) -- (leit);

    % ------------- anti-circularidade do envelope --------------------------
    \node[rot, text width=44mm, anchor=north] (anticirc) at (0.6,-7.5)
        {anti-circularidade do envelope: a aptidão do AG usa uma partição de
         aferição disjunta do teste; os indivíduos finais são reavaliados no
         conjunto de teste intocado, e é esse valor, tipicamente inferior,
         que os resultados reportam};
    \draw[seta, dashed] (anticirc.north) -- (ag.south);
\end{tikzpicture}
\end{document}
%   (no modo standalone, troque as \ref{...} por texto fixo ou defina-as)
% Requer apenas arrows.meta e positioning (já em packages.tex). Legível em
% P&B: origem = branco; protocolo = cinza claro; leitura = cinza médio.
% FREEZE respeitado: todos os números são os da própria Seção 3.6
% (47 tamanhos, 10 a 200.000, 30 repetições; 4 cenários, 10 tamanhos,
% 10 a 30.000), nenhum resultado afirmado — a leitura remete ao Cap. 4.
% ============================================================================
\begin{tikzpicture}[
    x=1cm, y=1cm,
    caixa/.style={draw, rounded corners=2pt, align=center, text width=42mm,
                  minimum height=11mm, inner sep=3pt, font=\footnotesize},
    protocolo/.style={caixa, fill=gray!8},
    leitura/.style={caixa, fill=gray!14},
    seta/.style={-{Stealth[length=2.4mm]}, thick},
    rot/.style={font=\scriptsize, align=center, inner sep=2pt}
]
    % ------------- três origens para o conjunto inicial --------------------
    % (o AG fica na linha de baixo, vizinho da sua nota de anti-circularidade:
    % com ele no meio, a seta tracejada atravessava a caixa do DRI-SL)
    \node[caixa, text width=44mm] (aleat) at (0.6,0)
        {\textbf{amostragem aleatória}\\(dispersão do sorteio):\\
         47 tamanhos de $L_0$, de 10 a 200.000; 30 repetições por tamanho};
    \node[caixa, text width=44mm] (drisl) at (0.6,-2.9)
        {\textbf{DRI-SL} (proposto):\\determinístico e sem acesso a
         rótulos; algoritmo na Seção~\ref{sec:metodo-drisl}};
    \node[caixa, text width=44mm] (ag) at (0.6,-5.8)
        {\textbf{algoritmo genético}\\(envelope: teto e piso):\\
         4 cenários (max/min de Acurácia\\e de Macro F1) em 10 tamanhos,\\
         de 10 a 30.000};

    % ------------- protocolo único e leitura -------------------------------
    \node[protocolo, text width=46mm] (proto) at (6.8,-2.9)
        {protocolo comum a toda origem:\\o classificador PVBin é treinado\\
         no $L_0$; Acurácia e Macro F1\\medidas no conjunto de teste};
    \node[leitura, text width=46mm] (leit) at (6.8,-6.4)
        {leitura nos resultados:\\
         sensibilidade (Seção~\ref{sec:res-l0-sens}),\\
         envelope (Seção~\ref{sec:res-l0-ag}),\\
         \mbox{DRI-SL} \textit{versus} aleatório e\\
         envelope (Seção~\ref{sec:res-l0-drisl}), reexecução\\
         independente (Seção~\ref{sec:res-l0-replay})};

    \draw[seta] (aleat.east) -- (proto.150);
    \draw[seta] (drisl.east) -- (proto.west);
    \draw[seta] (ag.east) -- (proto.210);
    \draw[seta] (proto) -- (leit);

    % ------------- anti-circularidade do envelope --------------------------
    \node[rot, text width=44mm, anchor=north] (anticirc) at (0.6,-7.5)
        {anti-circularidade do envelope: a aptidão do AG usa uma partição de
         aferição disjunta do teste; os indivíduos finais são reavaliados no
         conjunto de teste intocado, e é esse valor, tipicamente inferior,
         que os resultados reportam};
    \draw[seta, dashed] (anticirc.north) -- (ag.south);
\end{tikzpicture}
\end{document}
%   (no modo standalone, troque as \ref{...} por texto fixo ou defina-as)
% Requer apenas arrows.meta e positioning (já em packages.tex). Legível em
% P&B: origem = branco; protocolo = cinza claro; leitura = cinza médio.
% FREEZE respeitado: todos os números são os da própria Seção 3.6
% (47 tamanhos, 10 a 200.000, 30 repetições; 4 cenários, 10 tamanhos,
% 10 a 30.000), nenhum resultado afirmado — a leitura remete ao Cap. 4.
% ============================================================================
\begin{tikzpicture}[
    x=1cm, y=1cm,
    caixa/.style={draw, rounded corners=2pt, align=center, text width=42mm,
                  minimum height=11mm, inner sep=3pt, font=\footnotesize},
    protocolo/.style={caixa, fill=gray!8},
    leitura/.style={caixa, fill=gray!14},
    seta/.style={-{Stealth[length=2.4mm]}, thick},
    rot/.style={font=\scriptsize, align=center, inner sep=2pt}
]
    % ------------- três origens para o conjunto inicial --------------------
    % (o AG fica na linha de baixo, vizinho da sua nota de anti-circularidade:
    % com ele no meio, a seta tracejada atravessava a caixa do DRI-SL)
    \node[caixa, text width=44mm] (aleat) at (0.6,0)
        {\textbf{amostragem aleatória}\\(dispersão do sorteio):\\
         47 tamanhos de $L_0$, de 10 a 200.000; 30 repetições por tamanho};
    \node[caixa, text width=44mm] (drisl) at (0.6,-2.9)
        {\textbf{DRI-SL} (proposto):\\determinístico e sem acesso a
         rótulos; algoritmo na Seção~\ref{sec:metodo-drisl}};
    \node[caixa, text width=44mm] (ag) at (0.6,-5.8)
        {\textbf{algoritmo genético}\\(envelope: teto e piso):\\
         4 cenários (max/min de Acurácia\\e de Macro F1) em 10 tamanhos,\\
         de 10 a 30.000};

    % ------------- protocolo único e leitura -------------------------------
    \node[protocolo, text width=46mm] (proto) at (6.8,-2.9)
        {protocolo comum a toda origem:\\o classificador PVBin é treinado\\
         no $L_0$; Acurácia e Macro F1\\medidas no conjunto de teste};
    \node[leitura, text width=46mm] (leit) at (6.8,-6.4)
        {leitura nos resultados:\\
         sensibilidade (Seção~\ref{sec:res-l0-sens}),\\
         envelope (Seção~\ref{sec:res-l0-ag}),\\
         \mbox{DRI-SL} \textit{versus} aleatório e\\
         envelope (Seção~\ref{sec:res-l0-drisl}), reexecução\\
         independente (Seção~\ref{sec:res-l0-replay})};

    \draw[seta] (aleat.east) -- (proto.150);
    \draw[seta] (drisl.east) -- (proto.west);
    \draw[seta] (ag.east) -- (proto.210);
    \draw[seta] (proto) -- (leit);

    % ------------- anti-circularidade do envelope --------------------------
    \node[rot, text width=44mm, anchor=north] (anticirc) at (0.6,-7.5)
        {anti-circularidade do envelope: a aptidão do AG usa uma partição de
         aferição disjunta do teste; os indivíduos finais são reavaliados no
         conjunto de teste intocado, e é esse valor, tipicamente inferior,
         que os resultados reportam};
    \draw[seta, dashed] (anticirc.north) -- (ag.south);
\end{tikzpicture}
\end{document}
%   (no modo standalone, troque as \ref{...} por texto fixo ou defina-as)
% Requer apenas arrows.meta e positioning (já em packages.tex). Legível em
% P&B: origem = branco; protocolo = cinza claro; leitura = cinza médio.
% FREEZE respeitado: todos os números são os da própria Seção 3.6
% (47 tamanhos, 10 a 200.000, 30 repetições; 4 cenários, 10 tamanhos,
% 10 a 30.000), nenhum resultado afirmado — a leitura remete ao Cap. 4.
% ============================================================================
\begin{tikzpicture}[
    x=1cm, y=1cm,
    caixa/.style={draw, rounded corners=2pt, align=center, text width=42mm,
                  minimum height=11mm, inner sep=3pt, font=\footnotesize},
    protocolo/.style={caixa, fill=gray!8},
    leitura/.style={caixa, fill=gray!14},
    seta/.style={-{Stealth[length=2.4mm]}, thick},
    rot/.style={font=\scriptsize, align=center, inner sep=2pt}
]
    % ------------- três origens para o conjunto inicial --------------------
    % (o AG fica na linha de baixo, vizinho da sua nota de anti-circularidade:
    % com ele no meio, a seta tracejada atravessava a caixa do DRI-SL)
    \node[caixa, text width=44mm] (aleat) at (0.6,0)
        {\textbf{amostragem aleatória}\\(dispersão do sorteio):\\
         47 tamanhos de $L_0$, de 10 a 200.000; 30 repetições por tamanho};
    \node[caixa, text width=44mm] (drisl) at (0.6,-2.9)
        {\textbf{DRI-SL} (proposto):\\determinístico e sem acesso a
         rótulos; algoritmo na Seção~\ref{sec:metodo-drisl}};
    \node[caixa, text width=44mm] (ag) at (0.6,-5.8)
        {\textbf{algoritmo genético}\\(envelope: teto e piso):\\
         4 cenários (max/min de Acurácia\\e de Macro F1) em 10 tamanhos,\\
         de 10 a 30.000};

    % ------------- protocolo único e leitura -------------------------------
    \node[protocolo, text width=46mm] (proto) at (6.8,-2.9)
        {protocolo comum a toda origem:\\o classificador PVBin é treinado\\
         no $L_0$; Acurácia e Macro F1\\medidas no conjunto de teste};
    \node[leitura, text width=46mm] (leit) at (6.8,-6.4)
        {leitura nos resultados:\\
         sensibilidade (Seção~\ref{sec:res-l0-sens}),\\
         envelope (Seção~\ref{sec:res-l0-ag}),\\
         \mbox{DRI-SL} \textit{versus} aleatório e\\
         envelope (Seção~\ref{sec:res-l0-drisl}), reexecução\\
         independente (Seção~\ref{sec:res-l0-replay})};

    \draw[seta] (aleat.east) -- (proto.150);
    \draw[seta] (drisl.east) -- (proto.west);
    \draw[seta] (ag.east) -- (proto.210);
    \draw[seta] (proto) -- (leit);

    % ------------- anti-circularidade do envelope --------------------------
    \node[rot, text width=44mm, anchor=north] (anticirc) at (0.6,-7.5)
        {anti-circularidade do envelope: a aptidão do AG usa uma partição de
         aferição disjunta do teste; os indivíduos finais são reavaliados no
         conjunto de teste intocado, e é esse valor, tipicamente inferior,
         que os resultados reportam};
    \draw[seta, dashed] (anticirc.north) -- (ag.south);
\end{tikzpicture}
